% content/c01-wahrscheinlichkeit/en.tex — Chapter 1: Probability for physicists
% (English edition)
%
% Same structure, IDs and display equations as de.tex, which leads.  Migration of the
% subsection "Probability" of the first lecture (latex/lectures/en/lec01.tex:93-216); see
% evidence/migrations/maps/vl01-map.json.  Definitions: registry D53 to D56.  Derivations:
% evidence/derivations/c01/DERIVATION.md, same IDs.  Framing: docs/plan/BRIEF.md:45-46
% (author-stipulated, quoted verbatim here); hazard H1.1 in proposition
% c01.prop.punktgewicht.

\chapter{Probability for physicists}\label{c01}

The numbers that quantum mechanics delivers as its results are \emph{probabilities}. How
are such quantities to be understood, and how do we work with them? The right framework for
that is probability theory, and we build it up in this chapter. Do not expect anything
conceptually new: I take an intuitive understanding of probabilities, distributions and
expectation values for granted. What matters is the \emph{notation} in which we shall later
formulate quantum mechanics, and a language. To it belongs a stipulation that I make for the
whole course:
% docs/plan/BRIEF.md:46, first part of the sentence, verbatim.  GAP-TJO: the German wording
% in de.tex is the chapter author's translation (G-0022).  The second part of the sentence
% (measurements as elements of the dual space) belongs to Chapter 2 and is not quoted here.
\begin{equation}\label{c01.eq.sprache}\claim{c01.cl.sprache}{author-stipulated}
  \begin{aligned}
    &\text{\tpcifenglish{Thus our presentation starts with probability theory,}{Unsere Darstellung beginnt deshalb mit der}} \\
    &\text{\tpcifenglish{with notation and framing}{Wahrscheinlichkeitstheorie, mit Notation und Rahmung,}} \\
    &\text{\tpcifenglish{adapted to physics: we use the language}{die der Physik angepasst sind: Wir benutzen die Sprache}} \\
    &\text{\tpcifenglish{of preparations (= states)}{der Präparationen (= Zustände)}} \\
    &\text{\tpcifenglish{instead of “probability measures”.}{statt der „Wahrscheinlichkeitsmaße“.}}
  \end{aligned}
\end{equation}
What mathematics calls a probability distribution is thus called a preparation here, and
what it calls a random variable is called a measurement. Table~\ref{c01.tab.sprache} at the
end of Section~\ref{c01.s2} sets the words side by side.

To model an experiment in terms of probability we need exactly three things: the set of
possible outcomes, the preparation and the measurement. All three must be fixed before we
can even begin to calculate. We take them in this order.

\section{Outcomes and preparations}\label{c01.s1}

The first step in modelling any probabilistic experiment is to assign a set to it. One
obtains it by listing all possible outcomes.

\tpcdef{D53}

In this section and the next two, $\Outcomes$ is a finite set.

\begin{example}[Die and coins]\label{c01.ex.wuerfel}
  For a die the outcome set consists of the possible numbers of spots:
  $\Outcomes = \{1, 2, 3, 4, 5, 6\}$. For a coin,
  $\Outcomes = \{\text{\tpcifenglish{heads}{Kopf}}, \text{\tpcifenglish{tails}{Zahl}}\}$,
  for two coins, with H for heads and T for tails,
  $\Outcomes = \{\text{HH}, \text{HT}, \text{TH}, \text{TT}\}$. “The number of spots is
  even” is the event $S = \{2, 4, 6\}$.
\end{example}

\begin{exercise}[Two dice]\label{c01.exc.zwei-wuerfel}\claim{c01.cl.zwei-wuerfel}{derived}
  Describe the outcome set for \emph{two} dice. How many elements does it have?
\end{exercise}
\begin{solution}
  An outcome is a pair $(x_1, x_2)$ of the number of spots $x_1$ of the first die and the
  number of spots $x_2$ of the second. The outcome set is the set of all such pairs with
  $x_1, x_2 \in \{1, \dots, 6\}$. For each of the $6$ values of $x_1$ there are $6$ values
  of $x_2$; so it has $6 \cdot 6 = 36$ elements
  (Figure~\ref{c01.fig.ereignis}).
\end{solution}

\begin{tpcfigure}{c01.fig.ereignis}
  \caption{Outcome sets and events. (a) One die: six outcomes; the event “even number of
    spots” is the set of hatched cells. (b) Two dice: every point is a pair $(x_1, x_2)$,
    $36$ outcomes in all. The event “the sum of the spots is $7$” consists of the six
    pairs marked by squares.}
\end{tpcfigure}

Step two is the preparation. How does one model a fair die, a fair coin or a loaded die?
In a probabilistic theory by giving a probability for every outcome.

\tpcdef{D54}

The two conditions on $\pw$ are the only ones. The first holds because a negative
probability cannot be interpreted; the second because the probability that some outcome
occurs is always one. For the fair coin,
$\pw(\text{\tpcifenglish{heads}{Kopf}}) = \pw(\text{\tpcifenglish{tails}{Zahl}}) = \tfrac12$.
In what follows we call $\pw$ a \emph{weight} for short.

\begin{example}[Fair and loaded die]\label{c01.ex.gezinkt}\claim{c01.cl.gezinkt}{derived}
  For a fair die, $\pw(x) = \tfrac16$ for all $x$. Other weights describe \emph{loaded}
  dice. In this chapter the loaded die is the weight
  \begin{equation}\label{c01.eq.gezinkt}
    \pw(1) = \tfrac{1}{12}, \qquad
    \pw(2) = \pw(3) = \pw(4) = \pw(5) = \tfrac{2}{12}, \qquad
    \pw(6) = \tfrac{3}{12} .
  \end{equation}
  The sum is $\tfrac{1 + 4\cdot 2 + 3}{12} = 1$. The event “even number of spots”,
  $S = \{2,4,6\}$, has probability
  $\Prob[S] = 3\cdot\tfrac16 = \tfrac12$ for the fair die and
  $\Prob[S] = \tfrac{2 + 2 + 3}{12} = \tfrac{7}{12}$ for the loaded one.
\end{example}

\begin{tpcfigure}{c01.fig.gewichte}
  \caption{Three weights on the outcome set $\{1, \dots, 6\}$ of a die. The height of a bar
    is $\pw(x)$; in each picture the sum of all heights is $1$. (a) Fair die. (b) The
    loaded die of equation~\eqref{c01.eq.gezinkt}. (c) A point weight: the whole weight
    lies on the outcome $3$. Such weights describe deterministic preparations
    (Proposition~\ref{c01.prop.punktgewicht}).}
\end{tpcfigure}

\begin{proposition}[Properties of the probability]\label{c01.prop.wahrscheinlichkeit}\claim[proof-included]{c01.cl.wahrscheinlichkeit}{derived}
  Let $\pw$ be a weight on a finite outcome set $\Outcomes$, and let
  $S, T \subset \Outcomes$ be events. Then:
  \begin{enumerate}
    \item $0 \leq \Prob[S] \leq 1$, $\Prob[\emptyset] = 0$ and $\Prob[\Outcomes] = 1$.
    \item If $S_1, \dots, S_k$ are pairwise disjoint, that is, if no two of them have an
      outcome in common, then
      \begin{equation}\label{c01.eq.additiv}
        \Prob[S_1 \cup \dots \cup S_k] = \Prob[S_1] + \dots + \Prob[S_k] .
      \end{equation}
    \item $\Prob[\Outcomes \setminus S] = 1 - \Prob[S]$.
    \item $\Prob[S \cup T] = \Prob[S] + \Prob[T] - \Prob[S \cap T]$.
  \end{enumerate}
\end{proposition}
\begin{proof}
  On~2: every outcome $x$ of the union lies in exactly one $S_j$. The sum of the
  $\pw(x)$ over the union can therefore be arranged according to the $S_j$, and that is the
  right-hand side. On~1: $\Prob[S]$ is a sum of numbers $\pw(x) \geq 0$, hence
  $\Prob[S] \geq 0$; the empty sum is $0$, and $\Prob[\Outcomes] = 1$ is the second
  condition on a weight. On~3: $S$ and $\Outcomes \setminus S$ are disjoint with
  union $\Outcomes$, so by~2 $\Prob[S] + \Prob[\Outcomes \setminus S] = 1$.
  From this $\Prob[S] \leq 1$ follows too. On~4: $S \cup T$ is the disjoint union of
  $S$ and $T \setminus S$, and $T$ is the disjoint union of $T \setminus S$ and
  $S \cap T$. By~2, $\Prob[S \cup T] = \Prob[S] + \Prob[T \setminus S]$ and
  $\Prob[T] = \Prob[T \setminus S] + \Prob[S \cap T]$. Substituting gives the claim.
\end{proof}

Disjointness in~\eqref{c01.eq.additiv} is essential. If the events overlap, the right-hand
side counts outcomes twice; part~4 shows for two events by how much.
A function that assigns numbers $\geq 0$ to sets and satisfies \eqref{c01.eq.additiv} is called
a \emph{measure} in mathematics, and a \emph{probability measure} if
$\Prob[\Outcomes] = 1$. The Core of this course needs no more measure theory; the words
are only meant to make the literature readable.

Preparations can be mixed. If $\pw$ and $q$ are weights and $\lambda$ is a number
between $0$ and $1$, one can prepare according to $\pw$ with probability $\lambda$ and
according to $q$ otherwise. That too is a preparation, with the weight
\begin{equation}\label{c01.eq.mischung}\claim{c01.cl.mischung}{derived}
  \lambda\, \pw + (1-\lambda)\, q,
  \qquad
  \bigl(\lambda\, \pw + (1-\lambda)\, q\bigr)(x) = \lambda\, \pw(x) + (1-\lambda)\, q(x) .
\end{equation}
This is again a weight: all values are $\geq 0$ because $\lambda$, $1-\lambda$,
$\pw(x)$ and $q(x)$ are, and the sum over $x$ is
$\lambda \cdot 1 + (1-\lambda)\cdot 1 = 1$.

\section{Measurements and expectation value}\label{c01.s2}

Chapter~0 split every experiment into three stages. In classical probability theory each
stage has a simple mathematical form
\srcref{tpc-script-v0.2}{lec01, box “Every experiment in three steps”}.

\begin{greybox}[The three stages, classically]
  \begin{equation}\label{c01.eq.stadien}\claim{c01.cl.messung}{source-backed}
    \begin{aligned}
      &\text{\tpcifenglish{1. Preparation:}{1. Präparation:}} && \text{\tpcifenglish{a weight}{ein Gewicht}}\ \pw \\
      &\text{\tpcifenglish{2. Run:}{2. Durchführung:}} && \text{\tpcifenglish{a map of weights to weights}{eine Abbildung von Gewichten auf Gewichte}} \\
      &\text{\tpcifenglish{3. Measurement:}{3. Messung:}} && \text{\tpcifenglish{a function}{eine Funktion}}\ f\colon \Outcomes \to \mathbb{R}
    \end{aligned}
  \end{equation}
  At the end of the measurement stands one of the possible outcomes. Which maps occur in
  the second stage is not needed in this chapter; in quantum mechanics this picture will
  come back again and again.
\end{greybox}

\tpcdef{D55}

Note how generous this is: we impose no condition on a measurement, practically every
function on $\Outcomes$ may count as a measurement. On the preparation, by contrast, we
impose two. The value $f(x)$ is the measured value obtained when one knows for certain that
the outcome $x$ has occurred. The expectation value connects preparation and measurement: it
is the sum of the measured values weighted with the probabilities.

\begin{example}[Number of spots]\label{c01.ex.augenzahl}\claim{c01.cl.augenzahl}{derived}
  For a die the number of spots $f(x) = x$ is a measurement. For the fair die
  \begin{equation}\label{c01.eq.augenzahl-fair}
    \Ex[f] = \tfrac16\,(1 + 2 + 3 + 4 + 5 + 6) = \tfrac{21}{6} = \tfrac72 ,
  \end{equation}
  for the loaded die of equation~\eqref{c01.eq.gezinkt}
  \begin{equation}\label{c01.eq.augenzahl-gezinkt}
    \Ex[f] = \tfrac{1}{12}\,(1\cdot 1 + 2\cdot 2 + 2\cdot 3 + 2\cdot 4 + 2\cdot 5 + 3\cdot 6)
           = \tfrac{47}{12} .
  \end{equation}
  Neither of the two values is a number of spots. The expectation value need not be a
  possible measured value. Nor is the number of spots the only measurement on a die: “Is
  the number of spots even?” and “Is the number of spots six?” are measurements too, with
  yes or no as the answer.
\end{example}

\begin{tpcfigure}{c01.fig.erwartung}
  \caption{The expectation value of the number of spots as a centre of mass. If one thinks
    of the weights $\pw(x)$ as masses at the places $x$ of a beam, then $\Ex[f]$ is the
    place at which the beam balances: (a) for the fair die in the middle, at
    $\tfrac72$; (b) for the loaded die of equation~\eqref{c01.eq.gezinkt} shifted to the
    right, at $\tfrac{47}{12}$.}
\end{tpcfigure}

\subsection{Relative frequencies}\label{c01.s2.1}

Where does the formula for the expectation value come from? It does not fall from the sky.
One repeats the experiment $N$
times with the same preparation, where $N$ is a positive integer, and records the outcomes
$x_1, \dots, x_N$. Let $N(x)$ be
the number of runs with the outcome $x$. The \emph{relative frequency} of $x$ is
\begin{equation}\label{c01.eq.haeufigkeit}\claim{c01.cl.haeufigkeit}{derived}
  h_N(x) = \frac{N(x)}{N},
  \qquad
  h_N(x) \geq 0,
  \qquad
  \sum_{x \in \Outcomes} h_N(x) = 1 .
\end{equation}
The sum is $1$ because each of the $N$ runs has exactly one outcome, so that
$\sum_x N(x) = N$. The relative frequencies of a series of runs are themselves a weight.
For the mean of the measured values of a measurement $f$ over the series,
\begin{equation}\label{c01.eq.stichprobenmittel}\claim{c01.cl.stichprobenmittel}{derived}
  \frac{1}{N} \sum_{i=1}^{N} f(x_i)
  = \frac{1}{N} \sum_{x \in \Outcomes} N(x)\, f(x)
  = \sum_{x \in \Outcomes} h_N(x)\, f(x) ,
\end{equation}
because in the sum on the left the term $f(x)$ occurs exactly $N(x)$ times. The mean of a
series of runs is thus the expectation value of $f$ in the weight $h_N$. This equation is
exact and holds for every series.

And how is $h_N$ related to $\pw$? If we carry out the experiment $N$ times independently
with the same preparation, then with many repetitions the relative frequency $h_N(x)$ is
likely to lie close to $\pw(x)$,
\begin{equation}\label{c01.eq.lesart}\claim{c01.cl.lesart}{source-backed}
  h_N(x) \approx \pw(x)
  \qquad \text{\tpcifenglish{for large}{für große}}\ N .
\end{equation}
In this model, for large $N$ the outcome $x$ is likely to occur about $N \pw(x)$ times;
the sum of the measured values is then about
$\sum_x N \pw(x) f(x)$, and after division by $N$ the expectation value stands there
\srcref{tpc-script-v0.2}{lec01, before the formula for the expectation value}. This reading
is an interpretation in the stated model: in a single series $h_N$ may differ from $\pw$.
For independent repetitions of the same preparation an upper bound on the probability of
large deviations can be proved\zoomref{c01.v2}.

\begin{tpcfigure}{c01.fig.haeufigkeit}
  \demo{c01.demo.haeufigkeit}
  \caption{A series of throws with the loaded die of equation~\eqref{c01.eq.gezinkt},
    generated by a computer program with a fixed seed. Hatched bars: relative frequencies
    $h_N(x)$ after the first (a) $12$, (b) $120$ and (c) $1200$ throws of the same series.
    Heavy horizontal marks: the weights $\pw(x)$. In (a) the bars lie far from the marks,
    in (c) close to them. The bars of one picture always have total height $1$, as in
    equation~\eqref{c01.eq.haeufigkeit}.}
\end{tpcfigure}

\subsection{Calculating with expectation values}\label{c01.s2.2}

Measurements are functions and can therefore be added, multiplied by numbers and
multiplied with one another: $(\alpha f + \gamma g)(x) = \alpha f(x) + \gamma g(x)$ and
$(fg)(x) = f(x)\,g(x)$. A number $c$ is read here as the constant measurement with the
value $c$. We write $f \leq g$ if $f(x) \leq g(x)$ for all $x \in \Outcomes$.

\begin{proposition}[Rules for the expectation value]\label{c01.prop.erwartung}\claim[proof-included]{c01.cl.erwartung}{derived}
  Let $\pw$ be a weight on a finite outcome set $\Outcomes$, let $f$ and $g$ be
  measurements and $\alpha$, $\gamma$, $c$ real numbers. Then:
  \begin{enumerate}
    \item Linearity: $\Ex[\alpha f + \gamma g] = \alpha\, \Ex[f] + \gamma\, \Ex[g]$.
    \item Normalisation: $\Ex[c] = c$.
    \item Monotonicity: $f \leq g$ implies $\Ex[f] \leq \Ex[g]$. In particular
      $\Ex[f] \geq 0$ if $f \geq 0$.
    \item The expectation value lies between the smallest and the largest measured value:
      $\min_x f(x) \leq \Ex[f] \leq \max_x f(x)$.
  \end{enumerate}
\end{proposition}
\begin{proof}
  On~1:
  \begin{equation}\label{c01.eq.linear}
    \sum_{x \in \Outcomes} \pw(x)\,\bigl(\alpha f(x) + \gamma g(x)\bigr)
    = \alpha \sum_{x \in \Outcomes} \pw(x) f(x) + \gamma \sum_{x \in \Outcomes} \pw(x) g(x) .
  \end{equation}
  On~2: $\sum_x \pw(x)\, c = c \sum_x \pw(x) = c$. On~3: by~1,
  $\Ex[g] - \Ex[f] = \Ex[g - f] = \sum_x \pw(x)\,(g(x) - f(x))$, and every term is a
  product of two numbers $\geq 0$. On~4: with $m = \min_x f(x)$ and $M = \max_x f(x)$,
  $m \leq f \leq M$; by~3 and~2, $m \leq \Ex[f] \leq M$ follows.
\end{proof}

\subsection{Yes-no measurements}\label{c01.s2.3}

Of the infinitely many possible measurements the most important are those that take only
the values $0$ and $1$. They are the model for detectors that indicate the presence of a
particle, such as a photon detector: the detector says click or no click. Quantum mechanics
hinges on these two-valued measurements.
Such a measurement $f$ is determined by the set $S = \{x \in \Outcomes \colon f(x) = 1\}$.
One writes
\begin{equation}\label{c01.eq.indikator-def}\claim{c01.cl.indikator-def}{source-backed}
  \ind{S}(x) =
  \begin{cases}
    1, & x \in S, \\
    0, & \text{\tpcifenglish{otherwise}{sonst}}
  \end{cases}
\end{equation}
and calls $\ind{S}$ the \emph{indicator function} of the event $S$. Function and set carry
exactly the same information. Hold on to this: indicator functions \emph{are} measurements,
not a separate object alongside them. The set $S$ may be arbitrary. Three limiting cases:
$S = \Outcomes$ gives the constant measurement $1$, $S = \emptyset$ the constant measurement
$0$, and for a single outcome $a$, $\ind{\{a\}}(x)$ equals $1$ for $x = a$ and $0$
otherwise; that is the Kronecker delta.

\begin{remark}
  Indicator functions turn up wherever probabilities occur. In gambling the most important
  measurement of all is: have I won or not? For a scratch card the outcome set is the set of
  all combinations of numbers that can appear. The measurement is $1$ for a win and $0$
  otherwise. Written out, this function is enormously complicated, for one has to test every
  combination for whether it wins; and yet it is only an indicator function. All the
  complexity sits in the set $S$.
\end{remark}

The expectation value of an indicator function is the probability of its event:
\begin{equation}\label{c01.eq.indikator}\claim{c01.cl.indikator}{derived}
  \Ex[\ind{S}] = \sum_{x \in \Outcomes} \pw(x)\, \ind{S}(x) = \sum_{x \in S} \pw(x) = \Prob[S] .
\end{equation}
Figure~\ref{c01.fig.indikator} shows the calculation for the loaded die. The expectation
value of an indicator function thus has a special form: it is a pure sum of probabilities.
Probabilities are special expectation values. Conversely $\Prob$ fixes all
expectation values, for with $S = \{x\}$ one has $\pw(x) = \Prob[\{x\}]$ and therefore
$\Ex[f] = \sum_{x} \Prob[\{x\}]\, f(x)$.

\begin{tpcfigure}{c01.fig.indikator}
  \caption{A yes-no measurement on the loaded die of equation~\eqref{c01.eq.gezinkt}.
    (a) The indicator function of the event $S = \{2,4,6\}$, dashed like all
    measurements. (b) The weights $\pw(x)$ (outlines) and the products $\pw(x)\,\ind{S}(x)$
    (hatched): the indicator function leaves only the bars over $S$. Their heights add
    up to $\Ex[\ind{S}] = \Prob[S] = \tfrac{7}{12}$.}
\end{tpcfigure}

One calculates with indicator functions as with the sets:
\begin{equation}\label{c01.eq.indikator-regeln}\claim{c01.cl.indikator-regeln}{derived}
  \ind{S}\,\ind{T} = \ind{S \cap T},
  \qquad
  \ind{S}^2 = \ind{S},
  \qquad
  1 - \ind{S} = \ind{\Outcomes \setminus S} .
\end{equation}
This is checked outcome by outcome: the product $\ind{S}(x)\,\ind{T}(x)$ is $1$ exactly
when both factors are $1$, that is, when $x$ lies in $S$ and in $T$. The second
equation is the case $T = S$. The third: $1 - \ind{S}(x)$ is $0$ for $x \in S$ and $1$
otherwise.

Every measurement can be put together from yes-no measurements. For a measured value $v$
let $\{f = v\}$ be the event $\{x \in \Outcomes \colon f(x) = v\}$, “the measurement gives
$v$”. Then
\begin{equation}\label{c01.eq.werteverteilung}\claim{c01.cl.werteverteilung}{derived}
  f = \sum_{v} v\, \ind{\{f = v\}},
  \qquad
  \Ex[f] = \sum_{v} v\, \Prob[\{f = v\}] ,
\end{equation}
where $v$ runs through the distinct values that $f$ takes. The first equation holds
because at every place $x$ exactly one indicator function equals one, the one with
$v = f(x)$. If $f(x) = 0$, all weighted terms vanish. The second follows from it by linearity and equation~\eqref{c01.eq.indikator}. The
expectation value is the mean of the measured values weighted with the probabilities.

A completely worked example shows the whole apparatus in miniature. In quantum mechanics we
shall calculate a great many expectation values; in the end that should come very easily to
you. It pays to do many such calculations yourself.

\begin{example}[An experiment with three outcomes]\label{c01.ex.drei-ergebnisse}\claim{c01.cl.drei-ergebnisse}{derived}
  The three outcomes of $\Outcomes = \{1, 2, 3\}$ might be, say, the number of photons that a
  detector counts. As the preparation we choose, arbitrarily but admissibly,
  $\pw(1) = \tfrac12$ and $\pw(2) = \pw(3) = \tfrac14$: all values are $\geq 0$, and
  $\tfrac12 + \tfrac14 + \tfrac14 = 1$. A set with three elements has $2^3 = 8$ subsets;
  each is an event and has its indicator function. Table~\ref{c01.tab.ereignisse} lists all
  eight with their probabilities. How large is the probability that no outcome occurs? One
  always occurs, so $\Prob[\emptyset] = 0$. The one-element sets have the values of $\pw$,
  the two-element sets follow from additivity~\eqref{c01.eq.additiv}, and
  $\Prob[\Outcomes] = 1$. With that the entire information about the experiment is listed,
  and along the way we have already computed eight expectation values:
  $\Prob[\{2,3\}] = \tfrac12$ \emph{is} the expectation value of the indicator function
  $\ind{\{2,3\}}$.

  Now a measurement that is not an indicator function: $f(1) = -1$, $f(2) = 0$,
  $f(3) = +1$. Its expectation value is
  \begin{equation}\label{c01.eq.drei-ergebnisse}
    \Ex[f] = \tfrac12\cdot(-1) + \tfrac14\cdot 0 + \tfrac14\cdot(+1) = -\tfrac14 .
  \end{equation}
  It is negative, and that is no problem. Only probabilities must be $\geq 0$; measured
  values and expectation values may be negative.
\end{example}

\begin{tpctable}{c01.tab.ereignisse}
  \caption{All eight events of the experiment of Example~\ref{c01.ex.drei-ergebnisse}, with
    the values of their indicator function at the points $1$, $2$, $3$ and their probability
    for $\pw = (\tfrac12, \tfrac14, \tfrac14)$. Every row is a yes-no measurement and its
    expectation value.}
  \begin{tabular}{lcccc}
    \toprule
    $S$ & $\ind{S}(1)$ & $\ind{S}(2)$ & $\ind{S}(3)$ & $\Prob[S] = \Ex[\ind{S}]$ \\
    \midrule
    $\emptyset$ & $0$ & $0$ & $0$ & $0$ \\
    $\{1\}$ & $1$ & $0$ & $0$ & $1/2$ \\
    $\{2\}$ & $0$ & $1$ & $0$ & $1/4$ \\
    $\{3\}$ & $0$ & $0$ & $1$ & $1/4$ \\
    $\{1, 2\}$ & $1$ & $1$ & $0$ & $3/4$ \\
    $\{1, 3\}$ & $1$ & $0$ & $1$ & $3/4$ \\
    $\{2, 3\}$ & $0$ & $1$ & $1$ & $1/2$ \\
    $\{1, 2, 3\}$ & $1$ & $1$ & $1$ & $1$ \\
    \bottomrule
  \end{tabular}
\end{tpctable}

Measured values $-1$, $0$, $+1$ are not an invention for the exercise. For angular momentum
we shall meet, in Chapter~10, states $\ket{\lqn, \mqn}$ on which the angular-momentum
component $\Lj{3}$ has the value $\hbar \mqn$, with
\begin{equation}\label{c01.eq.drehimpuls}\claim{c01.cl.drehimpuls}{source-backed}
  \mqn \in \{-\lqn, -\lqn + 1, \dots, \lqn\},
  \qquad
  \lqn = 1\colon \quad \mqn \in \{-1, 0, 1\} .
\end{equation}
In units of $\hbar$ these are, for $\lqn = 1$, exactly the measured values of $f$
\srcref{tpc-script-v0.2}{lec14, eigenvalues of angular momentum}.

\begin{exercise}[Three outcomes]\label{c01.exc.drei}\claim{c01.cl.drei}{derived}
  Let $\Outcomes = \{1, 2, 3\}$ and $\pw(x) = \tfrac13$ for all $x \in \Outcomes$.
  Compute $\Prob[S]$ for $S = \{1, 2\}$ and the expectation value of the measurement
  $f(x) = x^2$.
\end{exercise}
\begin{solution}
  $\Prob[S] = \pw(1) + \pw(2) = \tfrac23$ and
  $\Ex[f] = \tfrac13\,(1 + 4 + 9) = \tfrac{14}{3}$.
\end{solution}

\begin{tpctable}{c01.tab.sprache}
  \caption{The same things in two languages. On the left the words of probability theory,
    on the right those of physics, which this course uses.}
  \begin{tabular}{lll}
    \toprule
    Symbol & Probability theory & Physics \\
    \midrule
    $\Outcomes$ & sample space & outcome set \\
    $\pw$, $\Prob$ & probability distribution, measure & preparation \\
    $f$ & random variable & measurement \\
    $\ind{S}$ & indicator function of the event $S$ & yes-no measurement \\
    $\Ex[f]$ & expectation value & expectation value \\
    \bottomrule
  \end{tabular}
\end{tpctable}

\section{Variance and determinism}\label{c01.s3}

The expectation value says where the measured values lie on average. The variance says
how far they scatter around this value. In quantum mechanics it plays a large part; we shall
meet it again in the uncertainty relation.

\tpcdef{D56}

In the definition, $f - \Ex[f]$ is the measurement with the values $f(x) - \Ex[f]$, the
deviation from the expectation value, and $(f - \Ex[f])^2$ is its square.

\begin{proposition}[Properties of the variance]\label{c01.prop.varianz}\claim[proof-included]{c01.cl.varianz}{derived}
  Let $\pw$ be a weight on a finite outcome set $\Outcomes$, $f$ a measurement
  and $\alpha$, $\gamma$ real numbers. Then:
  \begin{enumerate}
    \item $\Var[f] = \Ex[f^2] - \Ex[f]^2$.
    \item $\Var[f] \geq 0$.
    \item $\Var[f] = 0$ holds if and only if $f(x) = \Ex[f]$ for all outcomes $x$ with
      $\pw(x) > 0$.
    \item $\Var[\alpha f + \gamma] = \alpha^2\, \Var[f]$.
  \end{enumerate}
\end{proposition}
\begin{proof}
  We write $m = \Ex[f]$; this is a number. On~1:
  $(f - m)^2 = f^2 - 2 m f + m^2$. With linearity and normalisation from
  Proposition~\ref{c01.prop.erwartung},
  \begin{equation}\label{c01.eq.varianzformel}
    \Var[f] = \Ex[f^2] - 2m\, \Ex[f] + m^2 = \Ex[f^2] - 2m^2 + m^2 = \Ex[f^2] - m^2 .
  \end{equation}
  On~2 and~3: by definition $\Var[f] = \sum_x \pw(x)\,(f(x) - m)^2$. Every term
  is $\geq 0$, hence so is the sum. It is $0$ exactly when every term is $0$,
  and the term belonging to $x$ is $0$ exactly when $\pw(x) = 0$ or $f(x) = m$.
  On~4: by linearity $\Ex[\alpha f + \gamma] = \alpha m + \gamma$. Hence
  $\alpha f + \gamma - \Ex[\alpha f + \gamma] = \alpha (f - m)$, its square is
  $\alpha^2 (f - m)^2$, and the expectation value of that is $\alpha^2\, \Var[f]$.
\end{proof}

\begin{example}[Variance of the number of spots]\label{c01.ex.wuerfel-varianz}\claim{c01.cl.wuerfel-varianz}{derived}
  For the fair die and $f(x) = x$,
  $\Ex[f^2] = \tfrac16\,(1 + 4 + 9 + 16 + 25 + 36) = \tfrac{91}{6}$, and with
  equation~\eqref{c01.eq.augenzahl-fair}
  \begin{equation}\label{c01.eq.varianz-fair}
    \Var[f] = \tfrac{91}{6} - \bigl(\tfrac72\bigr)^2 = \tfrac{182 - 147}{12} = \tfrac{35}{12} .
  \end{equation}
  For the loaded die of equation~\eqref{c01.eq.gezinkt},
  \begin{equation*}
    \Ex[f^2] = \tfrac{1}{12}\,(1 + 2\cdot 4 + 2\cdot 9 + 2\cdot 16 + 2\cdot 25 + 3\cdot 36)
    = \tfrac{217}{12} ,
  \end{equation*}
  and with equation~\eqref{c01.eq.augenzahl-gezinkt}
  \begin{equation}\label{c01.eq.varianz-gezinkt}
    \Var[f] = \tfrac{217}{12} - \bigl(\tfrac{47}{12}\bigr)^2
            = \tfrac{2604 - 2209}{144} = \tfrac{395}{144} .
  \end{equation}
\end{example}

\begin{example}[Two outcomes]\label{c01.ex.muenze}\claim{c01.cl.muenze}{derived}
  Let $\Outcomes = \{0, 1\}$ with $\pw(1) = q$ and $\pw(0) = 1 - q$ for a number $q$
  between $0$ and $1$, and let $f(x) = x$. Then $\Ex[f] = q$, and because $f^2 = f$,
  also $\Ex[f^2] = q$. Hence
  \begin{equation}\label{c01.eq.zwei-ergebnisse}
    \Var[f] = q - q^2 = q\,(1 - q) = \tfrac14 - \bigl(q - \tfrac12\bigr)^2 .
  \end{equation}
  The variance is at most $\tfrac14$, and it equals $0$ exactly for $q = 0$ and $q = 1$:
  when the outcome is certain.
\end{example}

The same calculation holds for every yes-no measurement. By
equation~\eqref{c01.eq.indikator-regeln}, $\ind{S}^2 = \ind{S}$, hence
\begin{equation}\label{c01.eq.var-indikator}\claim{c01.cl.var-indikator}{derived}
  \Var[\ind{S}] = \Prob[S] - \Prob[S]^2 = \Prob[S]\,\bigl(1 - \Prob[S]\bigr) .
\end{equation}

\begin{tpcfigure}{c01.fig.varianz}
  \caption{(a), (b) Two weights with the same expectation value $\tfrac72$ of the number of
    spots $f(x) = x$ (dashed line) and different variance: the fair die of
    equation~\eqref{c01.eq.varianz-fair} and a weight that lies half on $3$ and half on
    $4$; for the latter $\Ex[f^2] = \tfrac{9 + 16}{2}$ and
    $\Var[f] = \tfrac{25}{2} - \tfrac{49}{4} = \tfrac14$. (c) The variance
    of a yes-no measurement as a function of $q = \Prob[S]$ by
    equation~\eqref{c01.eq.var-indikator}: it vanishes exactly at the two ends, where
    the answer is certain, and is largest at $q = \tfrac12$.}
\end{tpcfigure}

\subsection{Point weights}\label{c01.s3.1}

For an outcome $a \in \Outcomes$,
\begin{equation}\label{c01.eq.punktgewicht}\claim{c01.cl.punktgewicht-def}{source-backed}
  \pw_a(x) =
  \begin{cases}
    1, & x = a, \\
    0, & \text{\tpcifenglish{otherwise}{sonst}}
  \end{cases}
\end{equation}
is called the \emph{point weight} on $a$. It is a weight, and it is concentrated on one
point: the outcome $a$ is prepared with certainty. For every measurement $f$,
$\Ex[f] = f(a)$. Every weight is a mixture of point weights, and in exactly one way:
\begin{equation}\label{c01.eq.zerlegung}\claim{c01.cl.zerlegung}{derived}
  \pw = \sum_{a \in \Outcomes} \pw(a)\, \pw_a .
\end{equation}
To check this, one evaluates both sides at a place $x$: on the right only the term with
$a = x$ remains, and it is $\pw(x)$. The decomposition is unique because
$\pw = \sum_a c_a\, \pw_a$ gives $c_x = \pw(x)$ at once on evaluating at the place $x$.

\begin{proposition}[Determinism is a special case]\label{c01.prop.punktgewicht}\claim[proof-included]{c01.cl.punktgewicht}{derived}
  For a weight $\pw$ on a finite outcome set $\Outcomes$ the following are equivalent:
  \begin{enumerate}
    \item $\pw$ is a point weight.
    \item Every event has probability $0$ or $1$.
    \item Every measurement has variance $0$.
    \item $\pw$ is not a proper mixture: if $\pw = \lambda q + (1 - \lambda) r$ with
      weights $q$, $r$ and $0 < \lambda < 1$, then $q = r = \pw$.
  \end{enumerate}
\end{proposition}
\begin{proof}
  1 implies~2: for $\pw = \pw_a$, $\Prob[S] = \ind{S}(a)$, which is $0$ or $1$.

  2 implies~1: for every $x$, $\pw(x) = \Prob[\{x\}]$ is $0$ or $1$. The sum
  of all $\pw(x)$ is $1$; so exactly one $\pw(a)$ equals $1$ and all others are $0$.

  1 implies~3: for $\pw = \pw_a$, $\Ex[f] = f(a)$, and $a$ is the only outcome
  with $\pw(x) > 0$. By Proposition~\ref{c01.prop.varianz}, part~3, $\Var[f] = 0$.

  3 implies~2: by equation~\eqref{c01.eq.var-indikator},
  $\Prob[S]\,(1 - \Prob[S]) = \Var[\ind{S}] = 0$, hence $\Prob[S] = 0$ or $\Prob[S] = 1$.

  1 implies~4: let $\pw_a = \lambda q + (1 - \lambda) r$. For $x \neq a$,
  $0 = \lambda q(x) + (1 - \lambda) r(x)$. Both terms are $\geq 0$ and
  $\lambda$, $1 - \lambda$ are positive, hence $q(x) = r(x) = 0$. Because $q$ and $r$
  are weights, $q(a) = r(a) = 1$ remains.

  4 implies~1: we show that a weight that is not a point weight is a proper
  mixture. If $\pw$ is not a point weight, there are two distinct outcomes $a$
  and $b$ with $\pw(a) > 0$ and $\pw(b) > 0$. Let $\varepsilon$ be the smaller of the two
  numbers. Then
  \begin{equation}\label{c01.eq.echte-mischung}
    q = \pw + \varepsilon\,(\pw_a - \pw_b),
    \qquad
    r = \pw - \varepsilon\,(\pw_a - \pw_b)
  \end{equation}
  are weights: the sum of the values is $1 + \varepsilon\,(1 - 1) = 1$ in each case, and
  only $q(b) = \pw(b) - \varepsilon$ and $r(a) = \pw(a) - \varepsilon$ could become
  negative, but by the choice of $\varepsilon$ they are $\geq 0$. One has $q \neq r$ because
  $q(a) - r(a) = 2\varepsilon > 0$, and $\pw = \tfrac12 q + \tfrac12 r$.
\end{proof}

\begin{tpcfigure}{c01.fig.mischung}
  \caption{The construction of equation~\eqref{c01.eq.echte-mischung} for the loaded
    die $\pw$ of equation~\eqref{c01.eq.gezinkt}, with $a = 6$, $b = 1$ and
    $\varepsilon = \tfrac{1}{12}$. Starting from $\pw$ (left), the weight $\tfrac{1}{12}$ is
    moved once from $1$ to $6$ (middle, $q$) and once from $6$ to $1$ (right, $r$); the
    changed bars are hatched. Here $r$ is the fair die. The loaded die is the mixture
    $\pw = \tfrac12 q + \tfrac12 r$ and therefore not a point weight.}
\end{tpcfigure}

In the finite model we call a preparation deterministic if its weight is a point weight.
Then the outcome of every measurement is fixed: $f$ gives the value $f(a)$ with probability
$1$. This is what the definite predictions look like of which Chapter~0 spoke, and this
is how they fit into a probabilistic theory: as preparations that are not proper mixtures.
% GAP-TJO: reading point weights as the example of "deterministic as a subset of
% probabilistic" (BRIEF.md:45) is the chapter author's (G-0023 d).

\begin{merksatz}[Key statement]
  A weight is a point weight if and only if every measurement has variance $0$, and if and
  only if it is not a proper mixture of other weights. Every other weight is in exactly
  one way a mixture of point weights.
\end{merksatz}

Chapter~2 calls states that are not proper mixtures \emph{pure}. The pure states of the
finite classical model are thus the point weights.
Equation~\eqref{c01.eq.zerlegung} decomposes every weight uniquely into point weights.
The same weight can also be produced by preparing, with probability $\pw(a)$, the outcome
$a$ with certainty.
Geometrically the weights on $n$ outcomes form a body with $n$ corners, the point weights:
a segment, a triangle, a tetrahedron\zoomref{c01.v1}.
% GAP-TJO: reading the decomposition as missing information (taken out of the prose after
% REVIEW-lec1) is the chapter author's link of
% c01.eq.zerlegung with lec01.tex:83 (G-0023 d).

\begin{vertiefung}[id=c01.v1,
    titel={The set of all preparations is a simplex},
    wozu={The Core shows that every weight is uniquely mixed from point weights. Here this
      becomes a picture: for $n$ outcomes the weights form a segment, a triangle, a
      tetrahedron and so on, and the point weights are the corners.},
    voraussetzung={Section~\ref{c01.s3.1}; vectors in $\mathbb{R}^n$. The rest of this
      chapter does not need the block; Chapter~2 takes up the picture.},
    zurueck={The weights on $n$ outcomes form a body with $n$ corners, the point weights;
      every point in it is mixed from the corners in exactly one way.}]
  Let $\Outcomes = \{1, \dots, n\}$. A weight $\pw$ is then the list
  $(\pw(1), \dots, \pw(n))$ of its values, a vector in $\mathbb{R}^n$. The point weight
  $\pw_a$ is the vector with a $1$ at the place $a$ and zeros elsewhere.

  \begin{proposition}[The weights form a simplex]\label{c01.prop.simplex}\claim[proof-included]{c01.cl.simplex}{derived}
    The set of all weights on $\{1, \dots, n\}$ is the set of all mixtures
    $\lambda_1 \pw_1 + \dots + \lambda_n \pw_n$ of the $n$ point weights with
    $\lambda_a \geq 0$ and $\lambda_1 + \dots + \lambda_n = 1$. The point weights are
    affinely independent: from
    \begin{equation}\label{c01.eq.affin-unabhaengig}
      \alpha_1 \pw_1 + \dots + \alpha_n \pw_n = 0
      \qquad \text{\tpcifenglish{and}{und}} \qquad
      \alpha_1 + \dots + \alpha_n = 0
    \end{equation}
    it follows that $\alpha_1 = \dots = \alpha_n = 0$.
  \end{proposition}
  \begin{proof}
    Every such mixture is a weight, for the same reasons as for
    equation~\eqref{c01.eq.mischung}. Conversely every weight is such a mixture by
    equation~\eqref{c01.eq.zerlegung}, with $\lambda_a = \pw(a)$. On affine
    independence: the left-hand side of the first equation in
    \eqref{c01.eq.affin-unabhaengig} is the vector $(\alpha_1, \dots, \alpha_n)$; if it is
    $0$, all $\alpha_a = 0$. The second equation is not needed for this.
  \end{proof}

  The convex hull of $n$ affinely independent points is called a simplex in Chapter~2, and
  a theory whose state space is a simplex is called classical there. For $n = 2$ the
  simplex is a segment, for $n = 3$ a triangle (Figure~\ref{c01.fig.simplex}), for
  $n = 4$ a tetrahedron. The numbers $\pw(1), \dots, \pw(n)$ are at once the
  probabilities of the outcomes and the mixing weights of the corners.

  \begin{tpcfigure}{c01.fig.simplex}
    \caption{The set of all weights. (a) Two outcomes: the segment between the point
      weights $\pw_1 = (1,0)$ and $\pw_2 = (0,1)$ in the plane of the coordinates
      $\pw(1)$, $\pw(2)$; marked is $\pw = (\tfrac14, \tfrac34)$. (b) Three
      outcomes: a triangle with the point weights as corners. Marked are the uniform
      weight $(\tfrac13, \tfrac13, \tfrac13)$ at the centroid and the weight
      $(\tfrac12, \tfrac13, \tfrac16)$, which lies closer to $\pw_1$. Filled points are
      point weights, open circles proper mixtures.}
  \end{tpcfigure}
\end{vertiefung}

\begin{vertiefung}[id=c01.v2,
    titel={How close are relative frequencies to the probabilities?},
    wozu={The Core restricts the statistical reading to independent repetitions of the
      same preparation. Here an upper bound on the probability of a deviation of the
      relative frequency by at least $\varepsilon > 0$ is proved; this bound decreases
      with the positive integer number of repetitions.},
    voraussetzung={Sections~\ref{c01.s2} and~\ref{c01.s3} up to
      equation~\eqref{c01.eq.var-indikator}. The rest of the course uses only the result
      in words.},
    zurueck={The probability that the relative frequency of an event after $N$ independent
      repetitions deviates from its probability by at least $\varepsilon$ is at most
      $1/(4N\varepsilon^2)$.}]
  The first step holds for every weight and every measurement.

  \begin{lemma}[Chebyshev's inequality]\label{c01.lem.tschebyscheff}\claim[proof-included]{c01.cl.tschebyscheff}{derived}
    Let $\pw$ be a weight on a finite outcome set $\Outcomes$, $f$ a measurement
    and $\varepsilon > 0$. Let $S$ be the event that $f$ deviates from its expectation
    value by at least $\varepsilon$,
    $S = \{x \in \Outcomes \colon \abs{f(x) - \Ex[f]} \geq \varepsilon\}$. Then
    \begin{equation}\label{c01.eq.tschebyscheff}
      \Prob[S] \leq \frac{\Var[f]}{\varepsilon^2} .
    \end{equation}
  \end{lemma}
  \begin{proof}
    For $x \in S$, $(f(x) - \Ex[f])^2 \geq \varepsilon^2$, and for all $x$,
    $(f(x) - \Ex[f])^2 \geq 0$. Hence $(f - \Ex[f])^2 \geq \varepsilon^2\, \ind{S}$ at
    every place. With monotonicity and linearity (Proposition~\ref{c01.prop.erwartung}) and
    equation~\eqref{c01.eq.indikator},
    $\Var[f] \geq \varepsilon^2\, \Ex[\ind{S}] = \varepsilon^2\, \Prob[S]$ follows.
  \end{proof}

  Now to the model. Let $N$ be a positive integer and $\Outcomes$ finite. An outcome
  of $N$ repetitions is a list $(x_1, \dots, x_N)$
  of outcomes from $\Outcomes$. We \emph{assume} that the repetitions are carried out with
  the same preparation $\pw$ and independently of one another, and describe this by the
  \emph{product weight}
  $\pw^{(N)}(x_1, \dots, x_N) = \pw(x_1)\cdots\pw(x_N)$ on the set of all such lists.
  Expectation value, probability and variance in this weight are written $\Ex_N$,
  $\Prob_N$ and $\Var_N$.

  \begin{lemma}[Product weight]\label{c01.lem.produkt}\claim[proof-included]{c01.cl.produkt}{derived}
    The product weight $\pw^{(N)}$ is a weight. If $g$ and $w$ are measurements on
    $\Outcomes$ and $i \neq j$ are two of the numbers $1, \dots, N$, then for the
    measurements $(x_1, \dots, x_N) \mapsto g(x_i)$ and
    $(x_1, \dots, x_N) \mapsto g(x_i)\,w(x_j)$ on the lists
    \begin{equation}\label{c01.eq.produkt}
      \Ex_N[g(x_i)] = \Ex[g],
      \qquad
      \Ex_N[g(x_i)\,w(x_j)] = \Ex[g]\;\Ex[w] .
    \end{equation}
  \end{lemma}
  \begin{proof}
    A sum over all lists of products $u_1(x_1)\cdots u_N(x_N)$ is the product of the
    single sums,
    \begin{equation*}
      \sum_{x_1, \dots, x_N} u_1(x_1)\cdots u_N(x_N)
      = \Bigl(\sum_{x_1} u_1(x_1)\Bigr)\cdots\Bigl(\sum_{x_N} u_N(x_N)\Bigr) ;
    \end{equation*}
    this is the right-hand side multiplied out. With $u_l = \pw$ for all $l$ this gives the sum
    $1$, and all values of $\pw^{(N)}$ are $\geq 0$. With $u_i = \pw\, g$ and $u_l = \pw$
    otherwise, the single sum belonging to $i$ is $\Ex[g]$ and all others are $1$. With
    $u_i = \pw\, g$, $u_j = \pw\, w$ and $u_l = \pw$ otherwise, the single sums belonging
    to $i$ and $j$ are $\Ex[g]$ and $\Ex[w]$, all others $1$.
  \end{proof}

  \begin{theorem}[Relative frequencies in independent repetitions]\label{c01.thm.haeufigkeit}\claim[proof-included]{c01.cl.gesetz}{derived}
    Let $\pw$ be a weight on a finite outcome set $\Outcomes$, $S$ an event
    and $N$ a positive integer. Let $H_N$ be the relative frequency of $S$ in a list,
    \begin{equation*}
      H_N(x_1, \dots, x_N) = \tfrac{1}{N}\bigl(\ind{S}(x_1) + \dots + \ind{S}(x_N)\bigr) .
    \end{equation*}
    Then in the product weight $\pw^{(N)}$
    \begin{equation}\label{c01.eq.gesetz-momente}
      \Ex_N[H_N] = \Prob[S],
      \qquad
      \Var_N[H_N] = \frac{\Prob[S]\,(1 - \Prob[S])}{N} ,
    \end{equation}
    and for every $\varepsilon > 0$
    \begin{equation}\label{c01.eq.gesetz}
      \Prob_N\bigl[\,\abs{H_N - \Prob[S]} \geq \varepsilon\,\bigr]
      \;\leq\; \frac{\Prob[S]\,(1 - \Prob[S])}{N \varepsilon^2}
      \;\leq\; \frac{1}{4 N \varepsilon^2} .
    \end{equation}
  \end{theorem}
  \begin{proof}
    By equation~\eqref{c01.eq.produkt} with $g = \ind{S}$,
    $\Ex_N[\ind{S}(x_i)] = \Prob[S]$ for every $i$, and by linearity
    $\Ex_N[H_N] = \tfrac{1}{N}\cdot N\, \Prob[S] = \Prob[S]$. For the variance we put
    $g = \ind{S} - \Prob[S]$; then $\Ex[g] = 0$ and $\Ex[g^2] = \Var[\ind{S}]$. One has
    $H_N - \Prob[S] = \tfrac{1}{N}\sum_i g(x_i)$ and
    \begin{equation}\label{c01.eq.gesetz-rechnung}
      \begin{aligned}
        \Var_N[H_N]
        &= \frac{1}{N^2} \sum_{i} \Ex_N[g(x_i)^2]
          + \frac{1}{N^2} \sum_{i \neq j} \Ex_N[g(x_i)\,g(x_j)] \\
        &= \frac{1}{N^2}\cdot N\, \Var[\ind{S}] + 0 .
      \end{aligned}
    \end{equation}
    By equation~\eqref{c01.eq.produkt} every mixed term equals
    $\Ex[g]\,\Ex[g]$, that is $0$. With equation~\eqref{c01.eq.var-indikator} this is the
    second statement in \eqref{c01.eq.gesetz-momente}. The first inequality in
    \eqref{c01.eq.gesetz} is Lemma~\ref{c01.lem.tschebyscheff} for the measurement $H_N$ in
    the weight $\pw^{(N)}$; the second follows from $q\,(1-q) \leq \tfrac14$
    (equation~\eqref{c01.eq.zwei-ergebnisse}).
  \end{proof}

  The bound falls like $1/N$. It does not say that $h_N$ tends to $\pw$ in every series
  of runs, and it says nothing about a single series. It says that series with a large
  deviation carry little weight in the model. It is moreover crude: the actual probability
  of the deviation can be much smaller than the bound.

  For a coin the weight of $H_N$ can be given in full. Let
  $\Outcomes = \{0, 1\}$, $\pw(1) = q$ with $0 \leq q \leq 1$ and $S = \{1\}$; let $N$ be a
  positive integer. We use the convention $t^0 = 1$ also for $t = 0$: a product with no
  factors is $1$. A list with exactly $k$ ones has
  the weight $q^k (1-q)^{N-k}$, and there are $\binom{N}{k}$ such lists, as many as there
  are ways of choosing the $k$ places of the ones among the $N$ places. Hence
  \begin{equation}\label{c01.eq.binomial}\claim{c01.cl.binomial}{derived}
    \Prob_N\bigl[H_N = \tfrac{k}{N}\bigr] = \binom{N}{k}\, q^k\, (1-q)^{N-k},
    \qquad k = 0, 1, \dots, N .
  \end{equation}
  Figure~\ref{c01.fig.konzentration} shows these probabilities for the fair coin.

  \begin{tpcfigure}{c01.fig.konzentration}
    \caption{The weight of the relative frequency $H_N$ of “tails” for a fair coin,
      $q = \tfrac12$, by equation~\eqref{c01.eq.binomial} for $N = 4$, $16$ and $64$
      throws. Horizontally the value $k/N$ of $H_N$; the height of a line is its
      probability, in each picture in units of the largest value. The bracket marks the
      values with $\abs{H_N - \tfrac12} < \tfrac14$. As $N$ grows the weight gathers in
      the middle.}
  \end{tpcfigure}
\end{vertiefung}

\section{Continuous outcomes}\label{c01.s4}

So far $\Outcomes$ was finite. For quantum mechanics, however,
$\Outcomes = \mathbb{R}$ is decisive: a particle in one dimension has a position,
and the possible positions are the real numbers. How to treat probabilities on arbitrary
sets properly is the business of measure theory. We get round it: we take the important
properties of the finite case and make new definitions from them, and only for the special
case that we actually need in the course. In this section we consider only
preparations on $\mathbb{R}$ that admit a density. In this density model probabilities
are obtained by integration, rather than by summing density values over points.
Such a preparation is described by a \emph{probability density}: a function $\pw$ on
$\mathbb{R}$ with
\begin{equation}\label{c01.eq.dichte}\claim{c01.cl.dichte}{source-backed}
  \begin{aligned}
    \pw(x) &\geq 0,
    &\int_{-\infty}^{\infty} \pw(x)\,\dd x &= 1, \\
    \Prob[[a,b]] &= \int_a^b \pw(x)\,\dd x,
    &\Ex[f] &= \int_{-\infty}^{\infty} \pw(x)\, f(x)\,\dd x .
  \end{aligned}
\end{equation}
The two conditions of the first row correspond to those on a weight. For real endpoints
$a \leq b$ the second row gives the interval probability. The expectation value of a
measurement $f$ is used here provided
$\int_{-\infty}^{\infty} \pw(x)\,\abs{f(x)}\,\dd x < \infty$. Compared with the finite
case, only the sum has been replaced by an integral. In the literature one also finds for the
expectation value the notation $\int \Prob(\dd x)\, f(x)$, with
$\Prob(\dd x) = \pw(x)\,\dd x$. I am not especially fond of it and shall hardly use it;
you should merely be able to recognise it.

\begin{merksatz}[A density is not a probability]
  The value $\pw(x)$ of a density is not a probability. Probabilities are
  areas under the curve, that is, integrals of $\pw$ over sets. A single point has
  probability $0$, and $\pw(x)$ may be larger than $1$.
\end{merksatz}

Both can be seen in the uniform density on the interval from $0$ to $\tfrac12$:
\begin{equation}\label{c01.eq.gleich}\claim{c01.cl.gleich}{derived}
  \begin{aligned}
    \pw(x) &=
    \begin{cases}
      2, & 0 \leq x \leq \tfrac12, \\
      0, & \text{\tpcifenglish{otherwise}{sonst}},
    \end{cases}
    &\int_{-\infty}^{\infty} \pw(x)\,\dd x &= 2\cdot\tfrac12 = 1, \\
    \Prob[[a,a]] &= \int_a^a \pw(x)\,\dd x = 0 .
  \end{aligned}
\end{equation}
The density has the value $2$ and is normalised all the same. The last equation holds for
every density and every point $a$: an integral over an interval of length $0$ is $0$.

\begin{tpcfigure}{c01.fig.dichte}
  \caption{(a) The density of the Gaussian distribution of Example~\ref{c01.ex.gauss}. The
    probability that the outcome lies between real endpoints $a \leq b$ is the hatched
    area and not the height of the curve. (b) The uniform density of
    equation~\eqref{c01.eq.gleich}: it has height $2$, and the area under it is $1$.}
\end{tpcfigure}

\begin{example}[Gaussian distribution]\label{c01.ex.gauss}\claim{c01.cl.gauss}{derived}
  The density
  \begin{equation}\label{c01.eq.gauss}
    \pw(x) = \frac{1}{\sqrt{2\pi}}\, e^{-\frac12 x^2}
  \end{equation}
  is called the \emph{normal} or \emph{Gaussian distribution}, the bell curve. It is really
  the most important example in all of probability theory. It is normalised: the area under
  the curve is $1$\zoomref{c01.v3}. We prepare a particle in this state and at the end look
  at its position. How probable is an outcome $x \geq 0$? Because $\pw(-x) = \pw(x)$, the
  areas to the left and to the right of the origin are equal; together they are $1$.
  Hence
  \begin{equation}\label{c01.eq.gauss-halb}
    \int_0^\infty \pw(x)\,\dd x = \int_{-\infty}^0 \pw(x)\,\dd x = \tfrac12 .
  \end{equation}
  Asked physically: is the particle to the left or to the right of the origin? Intuitively
  the result was clear beforehand. But be careful with the argument: it works \emph{only}
  for symmetric densities. For an asymmetric density such an integral is by no means
  obvious.
\end{example}

For measures on $\mathbb{R}$ it is often particularly easy to compute the expectation values
of the coordinate function $f(x) = x$ and its powers. These expectation values are called
\emph{moments}, provided the integrals
$\int_{-\infty}^{\infty} \abs{x}^n\,\pw(x)\,\dd x$ are finite:
\begin{equation}\label{c01.eq.momente}\claim{c01.cl.momente}{source-backed}
  M_n = \int_{-\infty}^{\infty} x^n\, \pw(x)\,\dd x ,
  \qquad n = 0, 1, 2, \dots
\end{equation}
The $n$-th moment is thus nothing but the expectation value of the measurement $x^n$.
Always $M_0 = 1$: this is the normalisation, geometrically the area under the curve. The
moments that matter most in practice are $M_1$ and $M_2$; $M_1$ is called the
\emph{mean}. The variance of the coordinate is, as in the finite case,
\begin{equation}\label{c01.eq.var-momente}\claim{c01.cl.var-momente}{derived}
  \begin{aligned}
    \Var[\Prob] &= \int_{-\infty}^{\infty} (x - M_1)^2\, \pw(x)\,\dd x \\
    &= M_2 - 2 M_1\cdot M_1 + M_1^2\cdot M_0 = M_2 - M_1^2 ,
  \end{aligned}
\end{equation}
provided $\int_{-\infty}^{\infty} \abs{x}\,\pw(x)\,\dd x$ and
$\int_{-\infty}^{\infty} x^2\,\pw(x)\,\dd x$ are finite. For quantum mechanics the variance matters more than
the second moment; it returns with the uncertainty relation. Mean and variance
carry the same information as $M_1$ and $M_2$.

The moments other than $M_0$ need not exist.

\begin{example}[A density without a mean]\label{c01.ex.kein-mittelwert}\claim{c01.cl.kein-mittelwert}{derived}
  Let $\pw(x) = 1/x^2$ for $x \geq 1$ and $\pw(x) = 0$ otherwise. This is a density:
  \begin{equation}\label{c01.eq.kein-mittelwert}
    \begin{aligned}
      \int_1^{b} \frac{\dd x}{x^2} &= 1 - \frac{1}{b} \;\longrightarrow\; 1, \\
      \int_1^{b} x\cdot\frac{\dd x}{x^2} &= \ln b \;\longrightarrow\; \infty
      \qquad (b \to \infty) .
    \end{aligned}
  \end{equation}
  The integral for $M_1$ diverges: this preparation has no mean, and still less a
  variance.
\end{example}

\begin{vertiefung}[id=c01.v3,
    titel={Normalisation and moments of the Gaussian distribution},
    wozu={Example~\ref{c01.ex.gauss} asserts that the Gaussian density is normalised. Here
      this is calculated, together with mean and variance.},
    voraussetzung={Section~\ref{c01.s4}; integrals in two variables, polar coordinates,
      integration by parts. The Core uses only the result.},
    zurueck={The Gaussian density $e^{-x^2/2}/\sqrt{2\pi}$ is normalised, its mean is $0$
      and its variance $1$.}]
  Let $I = \int_{-\infty}^{\infty} e^{-x^2/2}\,\dd x$. The integral exists because
  $e^{-x^2/2} \leq e^{-\abs{x}/2}$ for $\abs{x} \geq 1$. One computes $I^2$ as an
  integral over the plane and changes to polar coordinates $x = r\cos\varphi$,
  $y = r\sin\varphi$ with $\dd x\,\dd y = r\,\dd r\,\dd\varphi$:
  \begin{equation}\label{c01.eq.gauss-norm}\claim{c01.cl.gauss-norm}{derived}
    \begin{aligned}
      I^2 &= \int_{-\infty}^{\infty}\!\int_{-\infty}^{\infty} e^{-(x^2 + y^2)/2}\,\dd x\,\dd y
           = \int_0^{2\pi}\!\int_0^\infty e^{-r^2/2}\, r\,\dd r\,\dd\varphi \\
          &= 2\pi\,\Bigl[-e^{-r^2/2}\Bigr]_0^\infty = 2\pi .
    \end{aligned}
  \end{equation}
  Because $I > 0$, $I = \sqrt{2\pi}$ follows, and the density \eqref{c01.eq.gauss} has
  integral $1$.

  That the expected position is zero is intuitively clear; it is proved all the same. For
  the moments,
  \begin{equation}\label{c01.eq.gauss-momente}\claim{c01.cl.gauss-momente}{derived}
    \begin{aligned}
      M_1 &= 0, \\
      M_2 &= \frac{1}{\sqrt{2\pi}}\int_{-\infty}^{\infty} x\cdot x\,e^{-x^2/2}\,\dd x \\
          &= \frac{1}{\sqrt{2\pi}}\Bigl[-x\,e^{-x^2/2}\Bigr]_{-\infty}^{\infty}
            + \frac{1}{\sqrt{2\pi}}\int_{-\infty}^{\infty} e^{-x^2/2}\,\dd x
          = 0 + 1 .
    \end{aligned}
  \end{equation}
  $M_1$ vanishes because $x\,e^{-x^2/2}$ is an odd function and its integral
  over $x \geq 0$ exists: $\int_0^\infty x\,e^{-x^2/2}\,\dd x = 1$. For $M_2$ we
  integrated by parts, with $x\,e^{-x^2/2}$ as the derivative of $-e^{-x^2/2}$; the
  boundary term vanishes because $x\,e^{-x^2/2}$ tends to $0$ for $x \to \pm\infty$. By
  equation~\eqref{c01.eq.var-momente} the variance is $M_2 - M_1^2 = 1$.
\end{vertiefung}

With this the language for this chapter's models is ready: preparations are weights on
finite outcome sets or the densities on $\mathbb{R}$ considered here, measurements are
functions, and the two together give expectation values when the integrals are finite. Chapter~2 lets go of the
assumption that preparations are weights on an outcome set and keeps only what one can do
with them: mix them and ask yes-no questions.
